
% 1
\begin{exercicio}
    Calcule os determinantes
    \begin{mathlist}[2]
        \item \left|\begin{array}{rr}
            -3 & -1 \\
            2 &  4
        \end{array} \right|
        \item \left|\begin{array}{rr}
            13 & 7 \\
            11 & 5
        \end{array} \right|
        \item \left|\begin{array}{rr}
            3 & 1 \\
            5 & 2
        \end{array} \right|
        \item \left|\begin{array}{cc}
            \log(a) &     \log(b) \\
            \sfrac{1}{2} & \sfrac{1}{4}
        \end{array} \right|
        \item \left|\begin{array}{cc}
            2m^2 & 2m^4-m \\
            m &  m^3-1
        \end{array} \right|
    \end{mathlist}
    \begin{answers}
        \begin{mathlist}[3]
            \item -10
            \item -12
            \item 1
            \item \log\frac{\sqrt[4]{a}}{\sqrt{b}}
            \item - m^{2}
        \end{mathlist}
    \end{answers}
\end{exercicio}

% 2
\begin{exercicio}
    Calcule os determinantes
    \begin{mathlist}[2]
        \item \left|\begin{array}{rrr}
            1 & 1 & 0 \\
            0 & 1 & 0 \\
            0 & 1 & 1
        \end{array} \right|
        \item \left|\begin{array}{rrr}
            1 & 3 &  2 \\
            -1 & 0 & -2 \\
            2 & 5 &  1
        \end{array} \right|
        \item \left|\begin{array}{rrr}
            -3 & 1 & -7 \\
            2 & 1 &  3 \\
            5 & 4 &  2
        \end{array} \right|
        \item \left|\begin{array}{rrr}
            9 & 7 & 11 \\
            -2 & 1 & 13 \\
            5 & 3 &  6
        \end{array} \right|
        \item \left|\begin{array}{rrr}
            0 & a & c \\
            -c & 0 & b \\
            a & b & 0
        \end{array} \right|
        \item \left|\begin{array}{rrr}
            2 & -1 & 0 \\
            m &  n & 2 \\
            3 &  5 & 4
        \end{array} \right|
    \end{mathlist}
  \begin{answers}
      \begin{mathlist}[3]
        \item 1
        \item -9
        \item 20
        \item 121
        \item b \left(a - c\right) \left(a + c\right)
        \item 4 m + 8 n - 26
    \end{mathlist}
  \end{answers}
\end{exercicio}


% 3
\begin{exercicio}
    Calcule os determinantes das matrizes usando o Teorema de Laplace
    \begin{mathlist}
        \item M = \left|\begin{array}{rrrr}
            3 & 4 &  2 &  1 \\
            5 & 0 & -1 & -2 \\
            0 & 0 &  4 &  0 \\
            -1 & 0 &  3 &  3
        \end{array} \right|
        \item M = \left|\begin{array}{rrrr}
            0 & a & b & 1 \\
            0 & 1 & 0 & 0 \\
            a & a & 0 & b \\
            1 & b & a & 0
        \end{array} \right|
        \item M = \left|\begin{array}{rrrr}
            1 & 3 & 2 & 0 \\
            3 & 1 & 0 & 2 \\
            2 & 3 & 0 & 1 \\
            0 & 2 & 1 & 3
        \end{array} \right|
        \item M = \left|\begin{array}{rrrrr}
            1 & 2 &  3 & 4 & 5 \\
            0 & a & -1 & 3 & 1 \\
            0 & 0 &  b & 2 & 3 \\
            0 & 0 &  0 & c & 2 \\
            0 & 0 &  0 & 0 & d
        \end{array} \right|
        \item M = \left|\begin{array}{cccccc}
            x & 0 & 0 & 0 & 0 & 0 \\ 
            a & y & 0 & 0 & 0 & 0 \\ 
            l & p & z & 0 & 0 & 0 \\ 
            m & n & p & x & 0 & 0 \\ 
            b & c & d & e & y & 0 \\ 
            a & b & c & d & e & z
        \end{array} \right|
    \end{mathlist}        
  \begin{answers}
      \begin{mathlist}[3]
        \item -208
        \item a^{2} + b^{2}
        \item 48
        \item a b c d
        \item x^{2} y^{2} z^{2}
    \end{mathlist}
  \end{answers}
\end{exercicio}

\begin{exercicio}
    Calcule o valor do determinante
    \( \left|\begin{array}{cccc}
    	1 & a & a & a \\
    	1 & 0 & a & a \\
    	1 & 1 & 0 & a \\
    	1 & 1 & 1 & 1
    \end{array} \right| \)
    \begin{answers}
        \(\qquad a^{2} \left( 1 - a \right) \)
    \end{answers}
\end{exercicio}

\begin{exercicio}
    Calcule o valor de $A$ sabendo que
    \[A = \left|\begin{array}{cccc}
    	1 & 2 & 4  &  8  \\
    	1 & 3 & 9  & 27  \\
    	1 & 4 & 16 & 64  \\
    	1 & 5 & 25 & 125
    \end{array} \right|\]
    \begin{answers}
          \(\qquad A = 12\)
    \end{answers}

\end{exercicio}
